\documentclass[11pt,a4paper]{article}

\usepackage[utf8]{inputenc}
\usepackage[T1]{fontenc}
\usepackage[provide=*,english]{babel}
\usepackage{lmodern}
\usepackage{microtype}
\usepackage{amsmath,amssymb,mathtools}
\usepackage{geometry}
\usepackage{xcolor}
\usepackage{enumitem}
\usepackage{booktabs}
\usepackage{fancyhdr}
\usepackage[hidelinks]{hyperref}

\geometry{left=2.3cm,right=2.3cm,top=2.2cm,bottom=2.25cm,headheight=14pt}
\setlength{\parindent}{0pt}
\setlength{\parskip}{0.38em}
\setlist[itemize]{topsep=0.25em,itemsep=0.18em,leftmargin=1.5em}
\allowdisplaybreaks

\definecolor{blue}{RGB}{34,73,110}
\definecolor{gray}{RGB}{90,96,102}

\pagestyle{fancy}
\fancyhf{}
\fancyhead[L]{\small\color{gray}GC2D: projected BM4}
\fancyhead[R]{\small\color{gray}Implicit 2: simultaneous solve}
\fancyfoot[C]{\thepage}
\renewcommand{\headrulewidth}{0.3pt}

\newcommand{\R}{\mathbb{R}}
\newcommand{\Id}{I}
\newcommand{\Y}{\mathcal{Y}}
\newcommand{\M}{\mathcal{M}}
\newcommand{\PhiBM}{\Phi^{\mathrm{BM4}}_{h,t_n}}
\newcommand{\norm}[1]{\left\lVert #1\right\rVert}

\begin{document}

\begin{center}
  {\LARGE\bfseries\color{blue}BM4 Implicit 2 with Hairer's Symmetric Projection}\par
  \vspace{0.35em}
  {\large Simultaneous Newton formulation for output and multiplier}\par
  \vspace{0.45em}
  {\small Schur-complement equivalent to the reduced projected BM4 method}
\end{center}

\vspace{0.5em}

\section{Purpose and relation to BM4 implicit 1}

\texttt{BM4Implicit2} computes the same exact projected physical map as
\texttt{BM4Implicit1}, but retains the duplicated output state as a Newton
unknown.  For one guiding-centre particle, implicit 1 solves only for a
two-component multiplier, whereas implicit 2 solves simultaneously for four
duplicated output components and two multiplier components.  Its Newton system
is therefore $6\times6$.

Both methods surround one complete twelve-stage BM4 cycle with Hairer's
symmetric projection.  Neither method projects after the individual BM4 stages;
that different operation belongs to \texttt{ProjectedBM4Composition}.

\section{Doubled state and projection operators}

For the physical guiding-centre state $z=(X,Y)\in\R^2$, define
\[
 \Y=(u,v)\in\R^4,
 \qquad
 \M=\{(u,v):u=v\}.
\]
The relevant constant matrices are
\[
 E=\begin{pmatrix}\Id_2\\\Id_2\end{pmatrix},
 \qquad
 G=(\Id_2\;-\Id_2),
 \qquad
 N=G^T=\begin{pmatrix}\Id_2\\-\Id_2\end{pmatrix}.
\]
Consequently,
\[
 GE=0,
 \qquad
 GN=2\Id_2.
\]
The input $Ez_n$ lies on the physical diagonal and $N\mu=(\mu,-\mu)$
is normal to it.

\section{The BM4 base map}

Let $D_{s,t}$ denote the direct GC extended map and $D^*_{s,t}$ its
adjoint.  The direct map updates the second copy, then the first copy, and then
applies the exact harmonic coupling.  The adjoint applies the coupling first and
then performs the triangular updates in reverse order.

The BM4 coefficients are
\[
\begin{aligned}
 a_1&= 0.0792036964311957,&
 a_2&= 0.1303114101821663,\\
 a_3&= 0.2228614958676077,&
 a_4&=-0.3667132690474257,\\
 a_5&= 0.3246481886897062,&
 a_6&= 0.1096884778767498.
\end{aligned}
\]
With
\[
 (b_1,\ldots,b_{12})
 =(a_1,a_2,a_3,a_4,a_5,a_6,a_6,a_5,a_4,a_3,a_2,a_1),
\]
the implementation alternates adjoint and direct stages.  Define
\[
 \tau_0=t_n,
 \qquad
 \tau_j=t_n+h\sum_{k=1}^j b_k.
\]
Then
\[
 \Y^{(j)}=
 \begin{cases}
  D^*_{b_jh,\,\tau_{j-1}}(\Y^{(j-1)}),&j\text{ odd},\\
  D_{b_jh,\,\tau_j}(\Y^{(j-1)}),&j\text{ even},
 \end{cases}
\]
and
\[
 \boxed{\PhiBM(\Y^{(0)})=\Y^{(12)}.}
\]
The sum of the twelve coefficients is one.  Their palindromic order and the
direct/adjoint pairing define a symmetric fourth-order map, including the
negative fourth coefficient.

\section{Simultaneous projection equations}

Start from the physical input $\Y_n=Ez_n$ and pre-correct it with the unknown
multiplier:
\[
 \widehat{\Y}_n(\mu)=Ez_n+N\mu.
\]
BM4 implicit 2 retains the diagonal output $\Y_{n+1}$ as a separate unknown and
imposes
\[
 d(\Y_{n+1},\mu)
 :=\Y_{n+1}-N\mu-\PhiBM(Ez_n+N\mu)=0, \tag{1}
\]
\[
 g(\Y_{n+1}):=G\Y_{n+1}=0. \tag{2}
\]
If $\Y_{n+1}=(y_u,y_v)$ and
$\PhiBM(Ez_n+N\mu)=(u_f(\mu),v_f(\mu))$, the residual is
\[
 d_u=y_u-\mu-u_f(\mu),
 \qquad
 d_v=y_v+\mu-v_f(\mu),
 \qquad
 g=y_u-y_v.
\]
Equations (1)--(2) express both the projected BM4 step and the final diagonal
constraint without eliminating any output variable.

\section{Simultaneous Newton matrix}

Set
\[
 J_{\mathrm{BM4}}
 =D\PhiBM(Ez_n+N\mu).
\]
Since $E$, $G$, and $N$ are constant,
\[
 D_{\Y_{n+1}}d=\Id_4,
 \qquad
 D_\mu d=-(\Id_4+J_{\mathrm{BM4}})N,
\]
\[
 D_{\Y_{n+1}}g=G,
 \qquad
 D_\mu g=0.
\]
At iteration $k$, the simultaneous Newton correction solves
\[
 \boxed{
 \begin{pmatrix}
  \Id_4&-(\Id_4+J_{\mathrm{BM4}})N\\[2pt]
  G&0
 \end{pmatrix}
 \begin{pmatrix}
  \Delta\Y_{n+1}\\
  \Delta\mu
 \end{pmatrix}
 =-
 \begin{pmatrix}
  d\\g
 \end{pmatrix}.
 } \tag{3}
\]
The updates are
\[
 \Y_{n+1}^{(k+1)}=\Y_{n+1}^{(k)}+\Delta\Y_{n+1},
 \qquad
 \mu^{(k+1)}=\mu^{(k)}+\Delta\mu.
\]
No inverse is formed explicitly.  For $N$ physical particles, the current code
assembles a dense $6N\times6N$ system.

\section{Exact SymPy expression for the BM4 Jacobian}

Write $W(z,t)=D_zf(z,t)$.  For $s_j=b_jh$, coupling frequency $\omega$,
and
\[
 R(\theta)=
 \begin{pmatrix}
  \cos\theta&-\sin\theta\\
  \sin\theta&\cos\theta
 \end{pmatrix},
\]
define the two coupling blocks
\[
 H_j=\frac{\Id_2+R(2\omega s_j)}{2},
 \qquad
 K_j=\frac{\Id_2-R(2\omega s_j)}{2}.
\]
The Jacobian of the exact linear coupling is therefore
\[
 C_j=
 \begin{pmatrix}
  H_j&K_j\\
  K_j&H_j
 \end{pmatrix}.
\]

For an even, direct stage $j$, let
\[
 \widetilde v_j
 =v^{(j-1)}+s_jf\bigl(u^{(j-1)},\tau_j\bigr),
\]
and set
\[
 U_j=W\bigl(u^{(j-1)},\tau_j\bigr),
 \qquad
 V_j=W(\widetilde v_j,\tau_j),
 \qquad
 L_j=\Id_2+s_j^2V_jU_j.
\]
SymPy multiplication of the two shears and the final coupling gives the exact
$4\times4$ direct-stage Jacobian
\[
 \boxed{
 \begin{aligned}
 \mathcal D_j
 &=C_j
 \begin{pmatrix}
  L_j&s_jV_j\\
  s_jU_j&\Id_2
 \end{pmatrix}\\[2mm]
 &=
 \begin{pmatrix}
  H_jL_j+s_jK_jU_j&s_jH_jV_j+K_j\\
  K_jL_j+s_jH_jU_j&s_jK_jV_j+H_j
 \end{pmatrix}.
 \end{aligned}
 }
\]

For an odd, adjoint stage, first write the coupled input as
\[
 \begin{pmatrix}\overline u_j\\\overline v_j\end{pmatrix}
 =C_j
 \begin{pmatrix}u^{(j-1)}\\v^{(j-1)}\end{pmatrix},
 \qquad
 u^{(j)}=\overline u_j+s_jf(\overline v_j,\tau_{j-1}).
\]
Now define
\[
 U_j=W(\overline v_j,\tau_{j-1}),
 \qquad
 V_j=W\bigl(u^{(j)},\tau_{j-1}\bigr),
 \qquad
 L_j=\Id_2+s_j^2V_jU_j.
\]
The exact adjoint-stage Jacobian is
\[
 \boxed{
 \begin{aligned}
 \mathcal A_j
 &=
 \begin{pmatrix}
  \Id_2&s_jU_j\\
  s_jV_j&L_j
 \end{pmatrix}C_j\\[2mm]
 &=
 \begin{pmatrix}
  H_j+s_jU_jK_j&K_j+s_jU_jH_j\\
  s_jV_jH_j+L_jK_j&s_jV_jK_j+L_jH_j
 \end{pmatrix}.
 \end{aligned}
 }
\]
The order of every $U_j$, $V_j$, $H_j$, and $K_j$ factor is essential: these
$2\times2$ matrices do not commute in general.

Let $\mathcal A_j[c]$ or $\mathcal D_j[c]$ mean that the corresponding stage
duration is $s_j=ch$.  Substitution of the palindromic BM4 coefficients and
application of the chain rule gives the requested explicit expression
\[
 \boxed{
 \begin{aligned}
 J_{\mathrm{BM4}}={}&
 \mathcal D_{12}[a_1]\mathcal A_{11}[a_2]
 \mathcal D_{10}[a_3]\mathcal A_9[a_4]
 \mathcal D_8[a_5]\mathcal A_7[a_6]\\
 &{}\times
 \mathcal D_6[a_6]\mathcal A_5[a_5]
 \mathcal D_4[a_4]\mathcal A_3[a_3]
 \mathcal D_2[a_2]\mathcal A_1[a_1].
 \end{aligned}
 }
\]
This factorization is exact, not a finite-difference approximation.  The
reproducible calculation in \texttt{bm4\_jacobian\_sympy.py} uses SymPy 1.14,
verifies both displayed stage identities, and preserves noncommutative block
order.  Even before the internal stage blocks are expanded, the fully expanded
twelve-factor product contains $2048$ ordered monomials in each final
$2\times2$ block, or $8192$ in total.  The factorized expression above is
therefore the explicit form suitable for analysis and implementation.

\subsection*{Centered-difference production evaluation}

The current implementation evaluates this exact derivative by centered
differences of the complete doubled map:
\[
 [J_{\mathrm{BM4}}]_{:,j}
 \approx
 \frac{\PhiBM(\widehat{\Y}+\varepsilon_je_j)
       -\PhiBM(\widehat{\Y}-\varepsilon_je_j)}{2\varepsilon_j},
\]
\[
 \varepsilon_j
 =\eta\max(1,|\widehat{Y}_j|),
 \qquad
 \eta=\sqrt[3]{\epsilon_{\mathrm{mach}}}.
\]
This choice reuses the exact production BM4 cycle and avoids maintaining a
second twelve-stage tangent implementation.  Its trade-off is the additional
map evaluations and a finite-difference floor in the Newton matrix.

\section{Schur complement and equivalence to implicit 1}

Eliminating $\Delta\Y_{n+1}$ from (3) gives
\[
 \mathcal K
 =G(\Id_4+J_{\mathrm{BM4}})N
 =2\Id_2+GJ_{\mathrm{BM4}}N. \tag{4}
\]
The reduced residual is
\[
 R(\mu)=G\PhiBM(Ez_n+N\mu)+2\mu.
\]
Therefore
\[
 R'(\mu)=\mathcal K.
\]
Moreover, applying $G$ to (1) gives
\[
 Gd=g-R(\mu).
\]
The Schur complement of the simultaneous system is consequently
\[
 \boxed{\mathcal K\Delta\mu=-R(\mu),}
\]
which is exactly the reduced Newton correction used by \texttt{BM4Implicit1}.
The formulations describe the same nonlinear root; they differ only in which
variables are retained in the linear system.

\section{Initial iterate and stopping rule}

The implementation starts from
\[
 \mu^{(0)}=0,
 \qquad
 \Y_{n+1}^{(0)}=\PhiBM(Ez_n).
\]
This makes $d=0$ initially, while the diagonal defect $g$ drives the first
correction.  For
\[
 s_z=\max(1,\norm{z_n}_\infty),
\]
the iteration stops when
\[
 \max\bigl(\norm d_\infty,\norm g_\infty\bigr)
 \leq \mathrm{atol}+\mathrm{rtol}\,s_z.
\]
At convergence, the neutral physical representative is
\[
 z_{n+1}=\frac{y_u+y_v}{2}.
\]

\section{Symmetry, symplecticity, and order}

At the exact root, equations (1)--(2) are the same symmetric projection
condition as in implicit 1.  If the root is locally unique, reversing the BM4
step and changing the sign of the multiplier yields
\[
 \Psi_{-h,t_{n+1}}\circ\Psi_{h,t_n}=\Id.
\]
The projected method retains fourth-order accuracy under the standard regularity
conditions.  Its projection multiplier is expected to scale as
\[
 \norm{\mu_h}_\infty=O(h^5).
\]
If the doubled BM4 cycle preserves the extended symplectic form, Hairer's exact
symmetric projection induces a symplectic physical map.  Finite residual
tolerances, centered-difference derivatives, and linear-solve round-off set the
observed structural-error floor.

\section{Reproducible comparison with BM4 implicit 1}

The numerical audit uses:
\begin{itemize}
 \item random potential amplitude $0.3$, maximum wave number $8$, grid
       $32\times32$, seed $27$, interpolation order $5$;
 \item $z_0=(1,1.2)$, $\rho=0.3$, coupling frequency $\omega=\pi/8$;
 \item time interval $[0,0.4]$, absolute and relative Newton tolerances
       $10^{-15}$, and at most $12$ iterations;
 \item centered-difference scale $\eta=\sqrt[3]{\epsilon_{\mathrm{mach}}}$.
\end{itemize}

\begin{center}
\small
\begin{tabular}{@{}ccccc@{}}
\toprule
$h$ & $\max|z^{(1)}-z^{(2)}|$ & max $\norm{\mu}_\infty$ & max residual & max it.\\
\midrule
$0.200$ & $6.66\times10^{-16}$ & $6.0004\times10^{-10}$ & $1.83\times10^{-15}$ & $1$\\
$0.100$ & $3.33\times10^{-16}$ & $2.1510\times10^{-11}$ & $1.31\times10^{-15}$ & $1$\\
$0.050$ & $2.22\times10^{-16}$ & $7.1747\times10^{-13}$ & $1.00\times10^{-15}$ & $1$\\
$0.025$ & $2.22\times10^{-16}$ & $2.3095\times10^{-14}$ & $1.44\times10^{-15}$ & $1$\\
\bottomrule
\end{tabular}
\end{center}

The superscripts $(1)$ and $(2)$ denote the reduced and simultaneous
formulations.  Their state differences remain at round-off.  Against a
$h=0.003125$ reference, the reduced formulation gives observed global orders
$4.013$, $4.002$, and $3.990$; the simultaneous formulation produces the same
states to the precision shown.  The multiplier orders $4.802$, $4.906$, and
$4.957$ approach the expected value five.

\section{Implementation correspondence}

The production implementation follows the mathematical blocks directly:
\begin{itemize}
 \item \texttt{BM4Implicit2} selects the simultaneous solver;
 \item \texttt{\_advance\_composition} evaluates the complete BM4 base map;
 \item \texttt{\_bm4\_map\_jacobian} evaluates $J_{\mathrm{BM4}}$;
 \item \texttt{\_solve\_simultaneous\_projected\_bm4\_step} assembles (1)--(3);
 \item nonlinear diagnostics record iterations, final residuals, and multiplier
       norms for every main integration step;
 \item the reusable comparison study runs both formulations on identical grids
       and reports aligned physical-flow diagnostics.
\end{itemize}

\end{document}
